\section{立法法}
\begin{enumerate}
  \item 3 - 4 分。交叉。绝对命题重点。
  \item 立法。\zd{创制、修改、补充、认可、废止}法律的活动。内容 \hz 效力。
  \item 广义，狭义 \yjt 享有国家立法权（中国 XX 法）的国家机关。区别，主体。
  \item 特点。立、执、司、守、监。最核心。
  \item 立法体制。立法权的分配体系。体制问题就是权利分配问题。
        \begin{enumerate}
          \item \zd{1 元 2 级，多层次多分支}。
                \begin{enumerate}
                  \item 1 元。在宪法下；2 级。中央和地方。
                  \item 多层次。不同效力等级的立法活动存在。效力上的多位阶。
                  \item 多分支。立法主体多样化。
                \end{enumerate}
          \item 宪法。
          \item 法律。\yi 基本法。人大、可授权人常 \er 普通法。人常。主席令。
                \begin{enumerate}
                  \item 法律保留。防范行政机关（国务院）。
                        \begin{enumerate}
                          \item 绝对保留。\zd{犯罪刑罚、政治人身、司法监察}只能人大、人常设定。下位可以具体规定。
                          \item 相对保留。\zd{授权设定}。5 年可延长，届满 6 月要报告。\zd{暂行}条例。
                        \end{enumerate}
                \end{enumerate}
          \item 行政法规。国务院。总理签署国务院令。监察法规 \ydy 行政法规。30 日报人常。
          \item 地方性法规。
                \begin{enumerate}
                  \item 2 级。\yi 省人大人常 \er 地级市人大人常。\zd{重大事项人大过}。
                  \item 协同立法。双方都得地级市以上。
                  \item 地级市的立法权限。\zd{城乡管}理、城乡\zd{建}设、\zd{生态}文明、\zd{历史}文化、\zd{基}层\zd{治}理。
                \end{enumerate}
          \item 规章。\yi 国务院部门 \er 地方（省、市）。
                \begin{enumerate}
                  \item 矛盾统统找国务院。
                  \item 于公于私俩限制。不能增加权力，减少职责；增加义务，减少权利。
                  \item 地级市的规章也是 5 项内容。老人老办法，新人新办法。
                  \item 人大过渡找政府。先让政府搞规章 \yjt 过渡期 2 年 \yjt 法规。同相对保留。
                \end{enumerate}
          \item 自治条例、单行条例。民族自治区、州、县（没有乡）人大。\zd{宪、民、基、专}门规定不能变。
          \item 3 区变通授权。经济特区、浦东新区、海南自贸港区。\yi 全人大人常授权 \er \ydy 法律 \san 冲突找人常。
          \item 授权。基本法、相对保留、地方变通授权（自下而上）、中央单独针对某些地方（自上而下，开恩）。
        \end{enumerate}
  \item 程序。比较恶心。
        \begin{enumerate}
          \item 原则。\zd{科学、合法、民主}。
                \begin{enumerate}
                  \item 科学。
                        \begin{enumerate}
                          \item 立法规划。5 年期、1 年期、专项。
                          \item 难事论证会（专业），大事听证会（公众）。应当开会。
                          \item 评估。事前（可行性）、事后（实施效果）。可以评估。
                          \item 明确性、可操作性。
                        \end{enumerate}
                  \item 合法。
                  \item 民主。草案公开 30 天，委员长会议可决定不公开。\zd{基层立法联系点}。
                \end{enumerate}
          \item 提案。同后面人大、人常的议案。
          \item 审议。草案人大提前 1 月给，人常 1 星期。组织研讨。
          \item 听取提案说明。
          \item 进入议程。
          \item 组织讨论。\yi 全体代表审议 \er 专委会。
          \item 审议报告。\zd{宪法和法律委员会}汇总各方意见，形成草案修改稿。\zd{合宪性问题要单独说明}。
          \item 表决稿。
          \item 审议程序。可以找相关人大代表\zd{列席}。一般讨论 3 次，紧急 1 回（港区国安法）。
          \item 提交表决。
                \begin{enumerate}
                  \item 撤回人大找主席团、人常找委员长会议同意。
                  \item 无法形成共识。主席团提出，大会决定，授权常委会干活。
                  \item 暂不交付表决 \yjt 单独表决试水失败。
                \end{enumerate}
          \item 终止。2 年没有推进。\zd{必要时候可以延期}，委员长会议决定。
          \item 表决。
                \begin{enumerate}
                  \item 一般。全人大。
                  \item 合并（多部同类项）、单独（分歧大，试水）。可以。只出现在常委的立法程序。
                \end{enumerate}
          \item 通过。
          \item 发布。
          \item 备案。但凡备案都是 30 提案，敢早不敢迟。变通备案 + 单独说明。
          \item 公布。
                \begin{enumerate}
                  \item \zd{机关公报}、报纸（全国、地方）、官方网站（人大、政府的法制信息网）。
                  \item 法律、法规扩大公布范围。提案说明、草案说明、审议结果报告、各方意见。
                \end{enumerate}
        \end{enumerate}
  \item 审查。最复杂的问题。合宪性审查。批准机关连带责任。
        \begin{enumerate}
          \item 事前。是否生效。只有这两个需要批准。
                \begin{enumerate}
                  \item 能力太次 \yjt 地级市法规，提到上一级人常。
                  \item 权力太大 \yjt 自治条例、单行条例，除了县提到省人常，区州到全人常。
                  \item 人常审合法性。
                \end{enumerate}
          \item 事后。处理结果。撤销 \hz 改正；审查：合理、合法。
                \begin{enumerate}
                  \item 领导关系。审查 2 个，处理 2 个。人大对人常、政府对部门、上级部门对下级部门。
                  \item 监督关系。合法 + 撤销。1 1。同级人常对同级政府；上一级人常对下一级人大、人常；授权对被授权。
                  \item 例外。3 看，看批准 \yjt 看关系 \yjt 看例外。
                        \begin{enumerate}
                          \item 地方人常监督变 2 1。
                          \item 自治区条例变 1 1。合法 + 撤。
                          \item 自治州、县统归全人常。主体例外。
                          \item 超纲。3 区法规报国务院、全人常。
                        \end{enumerate}
                  \item 被动事后审查。提请。法律以下，两央两高省人常提要求，其他提建议；规章国务院只提建议。
                \end{enumerate}
          \item 备案。找所有的上位法的制定机关，不要忘记国务院，人大不备案，规章最高只到国务院。
        \end{enumerate}
  \item 裁决。授权、立法、法律矛盾。全人常。
\end{enumerate}
